% !TeX root = ../papers/nc-p3.tex
\section{可換二平面を含む対称三平面の相手と二次曲線}
\label{sec:marked-commuting-plane-fibers}
\begin{theorem}\label{thm:marked-commuting-plane-fibers}
$G'\subset S$ を既約な対称跡零三平面とし、その自己交換子の包の次元を2とする。
調和二次鉛筆の基点 $B$ を含む滑らかな $E=V(f)$ を固定する。
$G'$ からの混合行列式が $f$ の非零倍となる相手三平面の全体は、
Grassmann多様体のPl\"ucker埋込みにおいて一本の滑らかな平面二次曲線の開部分である。
具体的に、$G'$ の可換二平面は $E$ 上の三次因子 $T$ を定め、$B=T+p$ と書ける。
$p\notin T$ なら二次曲線から二点を除き、$p\in T$ なら一点を除く。
全ての解で $g=5$ である。従って前節の固定計量例外も、\emph{対称側の部分空間}
の段階では被約な三角形および重なった三次因子の両方について扱える。
\end{theorem}
\begin{proof}
三交換子写像の核の直線が可換二平面 $H_0$ を一意に決める。
この平面に巡回行列 $Y$ がある。実際、二固有値で二次最小多項式の元があれば、
その二次元固有空間上で独立な可換元の制限はスカラーではなく、
一般結合の最小多項式は三次である。全元が平方零なら全てが対称な階数1以下であり、
二独立な対称階数一行列の和が階数2になることに反する。
よって $kI+H_0=k[Y]$ は三次元の可換環である。巡回ベクトルで $V$ をこの環と同定すれば、
自己随伴性より計量は非退化な積対合となる。この環のスペクトル $T$ は長さ3で、
一次式の制限を全環とする平面埋込みを持つ。すなわち環の単位を含む全切断で線型系を取り、
計量で双対平面を元の点平面に戻す。前節の積対合の議論により二次Veroneseの線型包は
$kI+H_0$ となる。$T$ は平面を張り、$I_T(2)$ は三次元でイデアル層を生成する。
実際一次で制限が全射なので $H^1(\mathcal I_T(1))=0$、また $H^2(\mathcal I_T)=0$ より
$\mathcal I_T$ は2正則である。$T\subset B\subset E$ がスキームとして従い、$B$ の長さは4である。

前節の $X:V\to G'$ の階数は今は2で、像は $H_0$ となる。実際 $H_0$ の基底を
最初の二元に取ると第三の係数は零である。$p=\mathbf P(\ker X)$ と置く。
$B$ の外で $A(v)$ の二次元核は $X(v),T_v$ がともに非零であればその包であり、
両列の外積と三次小行列の間の等式も成り立つ。両者の可能な共通零集合はそれぞれ
$B$ と $B\cup\{p\}$ の有限集合なので倍率は再び定数である。
特に $p\in B$ である。$p$ の外で $T_v\in G'$ なら、$G'$ 上の交代式の根基は
$X(v)$ が生成し、$T_v$ は $H_0$ 内に入る。この議論は単元成分を取って局所環上でも
成立するため、$B$ と $T$ は $p$ の外で等しく、$E$ の因子として $B=T+p$ である。

$R$ の像を $p$ で消える二一次式 $l,m$ で張り、$\ker R$ を二次式で書けば
$J=H^0(\mathcal I_T(2))$ である。消去二平面上 $R$ の階数が1以下では行列式が
零か可約である。階数2の基底では $l,m$ の基準持上げの三次式 $f_0$ に対し、
全ての三次式が一意な二平面座標 $(j_1,j_2)\in J^2$ を用いて
\[
                     f_0+l j_2-m j_1
\]
と書ける。この線型部分の核は $(l h,m h)$ で $l h,m h\in J$ を満たす一次式 $h$ からなる。
$p\notin T$ なら $l,m$ の制限は単位イデアルを生成し、$T$ 上 $h=0$ となるから核は零。
像は六次元の $H^0(\mathcal I_B(3))$ の全体である。従って任意の $zf$、$z\in k$、
ごとに二平面がただ一つあり、非零倍は $z\ne0$ の場合である。

$p\in T$ なら $T$ は滑らかな曲線の因子なので、$l,m$ が $T$ に生成するイデアルの
零化部分は $p$ の局所環のsocleの直線である。一次式の $T$ への制限は同型なので
核はちょうど一次元である。像の三次式は $p$ で特異で、$f$ 自身は滑らかなので
この特異条件が切る五次元部分空間をちょうど張る。$f_0$ は $p$ で滑らかである。
実際 $0\ne T_p\in H_0$ なので $X(v_*)=T_p$ となる $v_*$ を取れ、$v_*$ は $p$ の方向にない。
$df_0$ は $p$ で $m(v_*)l-l(v_*)m\ne0$ である。従って $zf-f_0$ が像に入る
$z$ は一意かつ非零で、二平面の全解は核のアフィン直線となる。

最後にこれらの直線を二列の外積でPl\"ucker座標に戻す。
$p\notin T$ では $z$ の一次係数の二列 $j_1,j_2$ は独立である。
従属なら $f=l j_2-m j_1$ は可約になってしまう。$p\in T$ では核の二列 $lh,mh$ も独立である。
$R$ の像への射影で定数外積は非零、一次係数は像と核の混合部分で非零、
二次係数は核の二次外積で非零となる。この三係数は独立なので閉包は滑らかな二次曲線である。
無限遠では両列とも核に入り、行列式は零となる。$p\notin T$ のときだけ有限の $z=0$ も除く。
$g=4$ では判別式に直線因子があるため滑らかな解の全ては $g=5$ である。
\path{computations/verify_commuting_plane_fibers.py}は剰余点が三角形に重なる前後での二種類の線型階数を再現する。
\end{proof}
